%% --------------------------------------------------------
\section{Перехідні стани та уявні частоти}
%% --------------------------------------------------------

\subsection{Класифікація стаціонарних точок}

Гессіан дає змогу класифікувати стаціонарні точки потенціальної поверхні енергії (ППЕ) за знаками власних значень:
\[
\lambda_i = \frac{\partial^2 E}{\partial q_i^2}.
\]

\begin{center}
    \begin{tblr}{
        colspec = {X[l,2] X[c,2] X[l,3]},
        row{1} = {font=\bfseries},
        hline{1,2,Z} = {solid},
            }
        Тип                & Уявних частот & Характеристика             \\
        Мінімум            & 0             & Локальний мінімум          \\
        Перехідний стан    & 1             & Сідлова точка 1-го порядку \\
        Сідло 2-го порядку & 2             & Нестабільна структура      \\
    \end{tblr}
\end{center}

\textbf{Уявна частота:} якщо \(\omega^2 < 0\), то \(\omega = i|\omega|\), і відповідний напрямок руху описує нестійку моду.

%% --------------------------------------------------------
\subsection{Інверсія аміаку (\ce{NH3}) як приклад}
%% --------------------------------------------------------

Молекула аміаку має тригранно-пірамідальну структуру симетрії $C_{3v}$: атом Нітрогену розташований над площиною трьох атомів Гідрогену. Дзер\-каль\-но-еквівалентна конфігурація виникає, коли атом N переходить на протилежний бік площини --- процес, відомий як \emph{парасолькова інверсія}.

Якщо ввести координату $h$, що задає відстань атома N від площини H$_3$, то потенціальна енергія має форму подвійного мінімуму:
\[
E(h) \approx E_0 + \tfrac{1}{2} k h^2 + \tfrac{1}{4}\lambda h^4, \qquad \lambda>0.
\]
Пірамідальні структури відповідають точкам $h=\pm h_0$ (локальні мінімуми), а планарна структура ($h=0$) є \textbf{перехідним станом}.



У рівноважній геометрії ($h = h_0$) усі власні значення гесіану додатні, що відповідає стабільному мінімуму:
\[
\lambda_i > 0 \quad \Rightarrow \quad \omega_i = \sqrt{\lambda_i/\mu_i} > 0.
\]

Для планарної конфігурації ($h = 0$) одне власне значення стає від’ємним, і відповідна частота набуває уявного значення:
\[
\omega = i|\omega|, \qquad \omega^2 < 0.
\]
Це свідчить про сідлову точку першого порядку — перехідний стан інверсії аміаку.


У класичному наближенні атом \ce{N} мав би подолати потенціальний бар’єр, однак квантова механіка допускає \emph{тунелювання}.  Енергетичне розщеплення між ними ($\Delta E \approx 0.79~\text{см}^{-1}$) відповідає частоті мікрохвильового переходу, що спостерігається експериментально в \emph{амонійному мазері}.

%% --------------------------------------------------------
\paragraph{Обчислення в \texttt{PySCF}}
%% --------------------------------------------------------


У файлі \inputcode{nh3_inversion.py} наведено повний приклад дослідження
інверсії молекули аміаку. Розрахунок демонструє, як із допомогою
\texttt{PySCF} можна чисельно побудувати профіль потенціальної енергії $E(h)$
та проаналізувати коливальні частоти у мінімумі й у перехідному стані.

Основні етапи скрипта:

\begin{enumerate}
    \item \textbf{Функція \inlinecode{energy\_func(h)}}
    Обчислює енергію системи при фіксованій висоті атома N над площиною трьох атомів H.
    Для кожного значення $h$ створюється новий об'єкт \inlinecode{gto.M},
    виконується \texttt{RHF}-розрахунок, і повертається значення енергії.

    \item \textbf{Функція \inlinecode{analyze\_point(h)}}
    Для заданої конфігурації обчислює:
    \begin{itemize}
        \item енергію $E(h)$;
        \item аналітичний градієнт \inlinecode{mf.Gradients().kernel()};
        \item гессіан \inlinecode{mf.Hessian().kernel()} та відповідні гармонічні частоти;
        \item кількість уявних частот (як ознаку перехідного стану).
    \end{itemize}
    Аналіз частот здійснюється через функцію \inlinecode{thermo.harmonic\_analysis()}.

    \item \textbf{Основний цикл аналізу}
    Задається набір висот $h$ (груба та уточнена сітка) і для кожної точки виконується \inlinecode{analyze\_point(h)}.
    Енергії переводяться у ккал/моль, градієнти контролюють локальність мінімуму.
    Для точок із малим градієнтом виводяться частоти коливань.

    \item \textbf{Пошук перехідного стану}
    Після проходу всіх точок програма вибирає:
    \begin{itemize}
        \item мінімум енергії $E_{\min}$ (стійкі геометрії);
        \item максимум енергії $E_{\text{TS}}$ поблизу $h = 0$;
    \end{itemize}
    Перевіряється, чи в кандидата на TS лише одна уявна частота,
    і якщо так, визначається бар’єр інверсії:
    \[
        \Delta E_{\text{бар}} = (E_{\text{TS}} - E_{\min}) \times 627.51 \text{ ккал/моль}.
    \]

    \item \textbf{Візуалізація профілю потенціальної поверхні}
    Побудова графіка $E(h)$ виконується за допомогою \texttt{matplotlib},
    результат зберігається у файл \texttt{nh3\_inversion.pdf}.
\end{enumerate}

Таким чином, програма дозволяє послідовно отримати:
\begin{itemize}
    \item рівноважні геометрії двох еквівалентних мінімумів (пірамідальних форм NH$_3$);
    \item перехідний стан із планарною структурою ($h = 0$);
    \item уявну частоту інверсії, що відповідає русі атома N через площину H$_3$;
    \item бар’єр інверсії, який кількісно характеризує тунельне розщеплення рівнів.
\end{itemize}